\chapter{Introduction}
\label{ch:intro}

This chapter introduces the problem that MemGate addresses: how a long-running
agent built on a large language model should decide what to remember. It
explains why larger context windows do not remove the need for that decision,
sets out the motivation and scope of the project, states the problem
precisely, lists the contributions, and describes how the rest of the report
is organised.

\section{Background of the project topic}

Large language models (LLMs) are now used as the core of \emph{agents}:
programs that plan, call tools, reflect on their own mistakes and carry on
working over many sessions
\cite{yao2023react,shinn2023reflexion,wang2024agentsurvey,xi2023agentsurvey}.
Coding assistants, research assistants, tutoring systems and personal
assistants are all expected to remember what was said days or weeks earlier.
An LLM, however, has no memory of its own between calls. Everything it knows
about the conversation has to be placed in its prompt, inside a fixed
\emph{context window} \cite{brown2020gpt3}.

The simplest way to give an agent memory is to append every message to a
growing prompt and send the whole history back to the model on every turn.
This works at first but does not scale. The cost of self-attention grows with
the square of the sequence length, and the key--value (KV) cache that the
model keeps grows with every token, so each reply becomes slower and more
expensive until the history no longer fits at all. Techniques that stretch the
window, such as position interpolation \cite{peng2024yarn}, delay that point
but do not remove it.

A larger window does not solve the quality problem either. Across frontier
models, accuracy drops by 30--50\% well before the advertised limit
\cite{chroma2025contextrot}; synthetic tests find effective context lengths
far shorter than the claimed ones \cite{hsieh2024ruler}; and how much a model
uses a passage depends on where it sits in a long input
\cite{liu2024lostmiddle}. These effects come from attention over long
sequences rather than from missing training data. What limits an agent is
therefore \emph{curation}, the choice of what goes into the prompt, rather than
raw capacity.

Several systems already manage agent memory. MemGPT \cite{packer2023memgpt}
pages a virtual context in and out of external storage, MemoryBank
\cite{zhong2024memorybank} applies a forgetting curve, Generative Agents
\cite{park2023generative} rank memories by recency, importance and relevance,
and Mem0 \cite{chhikara2025mem0} extracts and updates facts with an LLM.
Retrieval-augmented generation (RAG) \cite{lewis2020rag} stores everything and
retrieves the passages most similar to the query. Agent architectures now
treat memory as a module in its own right \cite{sumers2024coala}. These
designs show that tiered memory with forgetting works in practice. They are,
however, evaluated as complete systems, so it is not known how much each
individual decision inside them (what to keep, what to compress, what to
forget) is actually worth.

\section{Motivation and scope of the report}

\subsection{Motivation}
Poor curation has a measurable price. On the LongMemEval benchmark, a model
given oracle retrieval of the right evidence reaches about 92\% accuracy, while
the same model choosing its own evidence reaches about 58\%
\cite{wu2024longmemeval}. Given the right evidence the model answers
correctly, so the 34-point gap is a failure of \emph{selection}, not of
reasoning. A recent survey of agent memory puts it plainly: current systems
``do not solve the fundamental challenge: deciding what to remember and what
to forget'' \cite{zhang2024memorysurvey}. A rate--distortion view of these
decisions \cite{shannon1959rd,cover2006elements} has been proposed
\cite{colaco2026rd}, but as an analytical framework rather than a measured
system.

At Review~1 the panel asked whether a three-tier memory layer was still novel,
since commercial systems already ship one. The honest answer was that it is
not. The project therefore moved its contribution one level deeper, to the
\emph{decision policy} itself: building a layer in which each decision can be
isolated, swapped and measured, and then using it to find out which parts of
the familiar design actually earn their cost.

\subsection{Scope}
\textbf{In scope.}
\begin{itemize}
\item A memory controller that scores, compresses, stores, forgets and
  retrieves conversational memory in front of a frozen LLM, which is treated
  as a black box and never trained or fine-tuned.
\item A three-tier store (recent verbatim turns, compressed gists, durable
  facts with embeddings) held under a fixed storage budget.
\item Several interchangeable scoring policies: recency only, a hand-written
  heuristic, an LLM judge, and a learned scorer.
\item Evaluation on public long-conversation benchmarks (LoCoMo, REALTALK,
  BEAM, MemBench, PersonaChat, LongMemEval) with statistical tests that
  respect the clustering of questions by conversation.
\item A working chat application built on the same code, so that the memory
  decisions can be seen and tested by a user.
\end{itemize}

\textbf{Out of scope.} Training or fine-tuning the base LLM; image and audio
memory (the project is text only); production multi-user serving with
accounts and billing.

\section{Problem statement}

\begin{quote}
\emph{Given a conversation that grows without bound, a frozen LLM that can
only see a bounded context of $B$ tokens, and a memory store that may hold at
most $S$ tokens (or bytes), design and evaluate a memory decision policy that
chooses, for every message, whether to keep it verbatim, compress it, forget
it or recall it, so that the answers to later questions reach the model as
often as possible. Measure the contribution of each decision separately,
against controls that differ from it in exactly one factor.}
\end{quote}

The problem has two parts. The \emph{engineering} part is to build the layer
so that it is fast, never exceeds its budget, and works on real input. The
\emph{research} part is to find out, with proper statistics, whether
scored forgetting beats oldest-first forgetting, whether compressing beats
dropping, whether a learned scorer beats a hand-written one, and whether the
answer survives a change of corpus.

\section{Salient contribution}

The project makes the following contributions.
\begin{enumerate}
\item \textbf{A storage-budgeted evaluation protocol.} Bounding only the
  assembled context charges a store-everything baseline nothing for what it
  keeps. MemGate bounds the whole store, which changes the ranking of the
  policies (Section~\ref{sec:accounting}).
\item \textbf{Selection pays; compression does not.} At identical storage,
  fidelity and retrieval, scored forgetting gains $+9.37$ answer-recall points
  over FIFO, while compressing an evicted turn instead of dropping it costs
  $7.13$ points (Sections~\ref{sec:selection} and~\ref{sec:significance}).
\item \textbf{A test of the rate--distortion crossover} up to $182\times$
  compression, which is not reached (Section~\ref{sec:scaling}).
\item \textbf{Byte accounting.} The embedding index outweighs the text it
  indexes by about $12\times$ and reverses the policy ranking below roughly
  40\,KiB (Section~\ref{sec:bytes}).
\item \textbf{A negative result for LLM-judged salience.} A 0.5B judge loses
  to a hand-written heuristic by $5.87$ points (Section~\ref{sec:scorer}).
\item \textbf{A learned scorer} trained on generated and human conversations,
  with the direction in which it transfers and the one in which it does not
  (Section~\ref{sec:transfer}).
\item \textbf{End-task accuracy} through a fixed reader, showing how much of
  the recall advantage becomes correct answers (Section~\ref{sec:endtask}).
\item \textbf{Comparison with MemGPT-style and Mem0-style memory} rebuilt in
  the same harness (Section~\ref{sec:external}).
\item \textbf{Scale.} A test on chats of up to ten million tokens and a scorer
  retrained on 14.6 times the data (Sections~\ref{sec:scale}
  and~\ref{sec:retrain}).
\item \textbf{A working application.} A chat app that runs the shipped scorer,
  shows every keep, compress and forget decision with its reasons, accepts
  pasted text and uploaded documents, and was stress-tested on million-token
  chats and 200 simultaneous users (Sections~\ref{sec:app}
  and~\ref{sec:stress}).
\item \textbf{A research paper}, ``\paperTitle'', prepared in conference and
  journal versions.
\end{enumerate}

Several of these results are negative or corrective. We regard that as the
useful part: described only as a three-tier memory layer, MemGate would repeat
work others have done, whereas measured against controls it shows which parts
of that design are worth their cost.

\section{Organisation of report}

The rest of the report is organised as follows.
\textbf{Chapter~\ref{ch:lit}} surveys the literature on agent memory, prompt
and cache compression, retrieval and evaluation, compares the main systems,
and identifies the research gap.
\textbf{Chapter~\ref{ch:method}} describes the methodology and implementation:
the block diagram, the computing resources used, the scoring, routing and
budget algorithms with their equations and flowcharts, the training of the
learned scorer, the experimental design, and the chat application.
\textbf{Chapter~\ref{ch:results}} presents the results and their analysis,
with confidence intervals, together with the standards followed.
\textbf{Chapter~\ref{ch:adv}} discusses advantages, limitations and
applications, and \textbf{Chapter~\ref{ch:concl}} concludes and sets out the
future scope. The appendices contain the flowcharts of the software, data
sheets of the datasets and models, the list of software components, and the
papers prepared from this work.

\section*{Summary}
LLM agents need a policy that decides what to remember, because larger
context windows move the problem rather than remove it. Existing memory
systems are evaluated whole, so the worth of each decision is unknown. MemGate
is built to measure those decisions one at a time under a bounded store, and
the rest of this report describes how it was built and what it found.
